% ─── Chapter 7: Experiments & Iterative Development ───────────────
\chapter{Experiments \& Iterative Development}
\label{ch:experiments}

This chapter documents the iterative development process through which the
system was debugged, improved, and validated.  Each phase addresses a specific
problem, proposes a solution, and reports measured results.  The phases are
presented in chronological order.

\section{Phase 1: Baseline Establishment}

\textbf{Starting point.} The initial system used a 22-action native
\texttt{poke-env} space, Adler-32 hash-based categorical embeddings (20\,000
buckets), no positional encoding, a 4-layer post-norm transformer, and a
disabled LSTM.

\textbf{Observation.} The training pipeline ran end-to-end.  The easy
curriculum stage (random opponents) was solved at approximately 84\,000~steps,
reaching the 85\,\% win-rate threshold.  However, the medium stage plateaued at
roughly 50\,\% win rate by $\sim$100\,000~steps and did not improve further
even after 2.5~million steps.

\textbf{Key metrics:}
\begin{itemize}
  \item Easy stage: promoted at $\sim$84K steps (85\,\% win rate).
  \item Medium stage: plateau at $\sim$50\,\% win rate, persisted through 2.5M+ steps.
  \item Loss signature: high value loss, high entropy, low policy loss---indicating
        the value network was failing to predict returns while the policy was not
        making meaningful updates.
\end{itemize}

% TODO: Insert MLflow screenshot of baseline training curves showing:
% - Win rate over time (easy promotion then medium plateau)
% - Value loss, policy loss, entropy curves

\section{Phase 2: Validation Infrastructure}

\textbf{Problem.} Without a standardised evaluation protocol, checkpoint
quality could only be assessed by inspecting training-time metrics, which are
noisy and may not reflect true policy strength.

\textbf{Solution.} Implemented a validation pipeline with multiple protocols:

\begin{center}
\small
\begin{tabularx}{\textwidth}{l X}
  \toprule
  \textbf{Protocol} & \textbf{Method} \\
  \midrule
  Smoke        & 10 episodes vs.\ random opponent (quick sanity check). \\
  Fixed paired & 20 frozen teams $\times$ 2 opponents = 40 battles per run. \\
  Mirror       & Both sides use identical team. \\
  \bottomrule
\end{tabularx}
\end{center}

Additionally, 20~frozen Gen\,8 random-battle teams were generated and stored
as JSON manifests (10~swapped pairs for fairness).

\textbf{First validation results} (checkpoint at step~51\,200):
\begin{itemize}
  \item Overall win rate: 37.5\,\%
  \item Vs.\ random: 75\,\%
  \item Vs.\ heuristic: 0\,\%
  \item Fallback events: 2.60~per battle (0.069~per step)---indicating frequent
        invalid-action recoveries.
\end{itemize}

A subsequent diagnostics run with the final checkpoint showed:
\begin{itemize}
  \item Overall: 50\,\%
  \item Vs.\ random: 95\,\%
  \item Vs.\ heuristic: 5\,\%
\end{itemize}

The sharp contrast between random and heuristic performance confirmed that the
agent had learned basic move selection but not strategic play.

\section{Phase 3: Vocabulary \& ID Migration}

\textbf{Problem.} Adler-32 hashing with modulo 20\,000 produced an estimated
7.5\,\% collision rate among $\sim$1\,800 unique entity strings.  Collisions
mean different Pok\'{e}mon species (or items or abilities) share the same
embedding, degrading the model's ability to distinguish between entities.

\textbf{Solution.} Replaced hash-based IDs with stable dictionary lookups:

\begin{itemize}
  \item Generated explicit vocabularies from Showdown data files.
  \item Mapped each entity to a unique integer index.
  \item Added support for display-form aliases (e.g.\ ``Shellos (East Sea)''
        $\to$ internal key).
\end{itemize}

\textbf{Result.} Vocabulary sizes: 1\,516~species, 584~items, 315~abilities.
Coverage audit confirmed \textbf{zero unknown entities} across the BDSP
dataset and validation manifests.  This change intentionally broke
compatibility with all previous checkpoints---new training runs are treated as
a new experiment family.

\section{Phase 4: Logging \& Metrics Standardisation}

\textbf{Problem.} Metric flow from environment to MLflow had gaps:
``unknown'' opponent metadata was not properly logged, and opponent-specific
metrics were sometimes missing from training runs.

\textbf{Solution.} Traced the full metric path from environment through
\texttt{EnvBridge} to MLflow.  Fixed root causes:
\begin{itemize}
  \item ``Unknown'' opponent metadata now correctly logged as a missing count.
  \item Opponent-specific keys (e.g.\
        \texttt{training/win\_rate\_vs\_random},
        \texttt{training/win\_rate\_vs\_heuristic}) now present in all runs.
  \item Defined a standard contract for high-signal metrics.
\end{itemize}

\textbf{Result.} All subsequent training runs produce complete, structured
metrics enabling reliable run comparison in MLflow.

\section{Phase 5: Embedding Audit \& Global Features}

\textbf{Problem.} The observation tensor lacked high-value global context
features that could help the model understand battle dynamics (e.g.\ which
opponent type it was facing, whether it was forced to switch).

\textbf{Solution.} Added to the global token (position~0) without changing
the tensor dimensions:
\begin{itemize}
  \item Opponent label categories (\texttt{opponent\_random},
        \texttt{opponent\_heuristic}, \texttt{opponent\_other}).
  \item Normalised battle turn, force-switch flag, trapped flag.
  \item Normalised available move and switch counts.
  \item Dynamax, Mega Evolution, and Z-Move availability flags.
\end{itemize}

Also discovered and fixed \textbf{broken weather encoding}: the battle state
passed weather as a Python dictionary rather than enum keys, so weather
information was never actually encoded into the observation.

\textbf{Result.} The model now receives richer context per step.  However,
some potentially high-value features remain missing (move PP, revealed vs.\
unrevealed opponent knowledge, speed control information).

\section{Phase 6: Position \& Role Embeddings}

\textbf{Problem.} The transformer received 13~tokens with no positional
information.  Self-attention is permutation-invariant, meaning the model
could not structurally distinguish between the active Pok\'{e}mon and bench
slots, or between own-team and opponent-team tokens.

\textbf{Solution.} Added two sets of learnable embeddings after input
projection:
\begin{itemize}
  \item \textbf{Position embeddings:} 13~vectors, one per token position.
  \item \textbf{Role embeddings:} 5~coarse roles (CLS, our active, our bench,
        opponent active, opponent bench), shared across positions within a role.
\end{itemize}

\textbf{Result.} The transformer now has explicit slot and side information
before self-attention.  Old checkpoints are incompatible due to new embedding
parameters---a fresh training run was required.

\section{Phase 7: Compressed Action Space}

\textbf{Problem.} The native \texttt{poke-env} action space has 22~actions
(6~switches + 4~regular moves + 12~gimmick variants).  This large branching
factor slows exploration, and the 12~gimmick variants are rarely all available
simultaneously.

\textbf{Solution.} Introduced a compressed 14-action space:

\begin{center}
\small
\begin{tabularx}{0.7\textwidth}{c l l}
  \toprule
  \textbf{Indices} & \textbf{Category} & \textbf{Mapping} \\
  \midrule
  0--3  & Regular moves & Move 1--4 \\
  4--7  & Gimmick moves & Auto: Mega $\to$ Z $\to$ Dynamax \\
  8--13 & Switches      & Party slot 1--6 \\
  \bottomrule
\end{tabularx}
\end{center}

Gimmick slots automatically resolve to the first legally available gimmick
type, so the agent only needs to learn ``use gimmick'' rather than ``use
specific gimmick variant.''

\textbf{Result.} Action mask changed from $(22,)$ to $(14,)$.  The reduced
branching factor should simplify exploration.  Incompatible with old
checkpoints.

\section{Phase 8: LSTM Memory Fix}

\textbf{Problem.} The LSTM path existed in model code but was non-functional
due to RLlib's connector path not supporting recurrent state propagation in
the production setup.

\textbf{Solution.} Rewrote the LSTM integration:
\begin{itemize}
  \item LSTM now runs even when no incoming recurrent state is provided
        (creates zero initial state).
  \item \texttt{compute\_values()} uses the same forward path as policy logit
        computation, ensuring consistent state handling.
  \item Single-step and sequence batches both pass through the recurrent path
        correctly.
  \item \texttt{get\_initial\_state()} returns unbatched state tensors of shape
        $[H]$ for compatibility with RLlib.
\end{itemize}

\textbf{Result.} Cross-turn memory is now architecturally supported.
End-to-end training validation with LSTM enabled remains a next step.

\section{Observation Pipeline Investigation}

\textbf{Problem.} Despite architectural improvements, the model trained on
Gen\,5 random battles showed 0\,\% win rate against heuristic opponents.

\textbf{Approach.} Systematic investigation using a pipeline validator and
debug logging to trace observation construction step-by-step.

\textbf{Three bugs discovered:}

\begin{enumerate}
  \item \textbf{Weather never encoded.} \texttt{battle.weather} was passed as
        a Python dictionary object rather than being iterated over its enum
        keys.  The weather feature was always zero.
  \item \textbf{Switch token--action index mismatch.} When a Pok\'{e}mon
        fainted, bench positions shifted but the switch-action mapping was not
        updated, causing the agent to switch to the wrong slot.
  \item \textbf{Action mask fallback contamination.} When the mask construction
        encountered an edge case (e.g.\ no legal actions), it set
        \texttt{mask[0] = 1.0} as a fallback, forcing the agent to always
        pick action~0.  This affected \textbf{14\,\% of training steps}.
\end{enumerate}

\textbf{After fixes:}
\begin{itemize}
  \item Win rate vs.\ random improved from 65\,\% to 85\,\%.
  \item Win rate vs.\ heuristic remained at $\sim$7.5\,\%.
  \item Value function showed negative explained variance ($-0.003$), suggesting
        fundamental reward-signal issues.
\end{itemize}

\begin{keyinsight}
Fixing observation bugs significantly improved performance against random
opponents but did not transfer to heuristic opponents.  This indicated that the
observation pipeline was no longer the bottleneck---the reward signal was.
\end{keyinsight}

% TODO: Insert MLflow comparison chart: before vs after observation bug fixes,
% showing win rate curves against random and heuristic opponents.

\section{Attention Collapse Analysis}

\textbf{Problem.} The 4-layer post-norm transformer showed concerning
attention patterns.  Was the model actually using its depth?

\textbf{Approach.} Extracted per-head, per-layer attention matrices from a
trained checkpoint (step~2\,199\,834).

\textbf{Results:}

\begin{center}
\small
\begin{tabularx}{\textwidth}{c l X}
  \toprule
  \textbf{Layer} & \textbf{Attention Pattern} & \textbf{Key Observation} \\
  \midrule
  Layer 0 & Structured &
    CLS$\to$our\_active: 25\,\%, CLS$\to$opp\_active: 12.5\,\%.
    Some opponent awareness. \\
  Layer 1 & Collapsed &
    Opponent attention drops to 3\,\% (12:1 ratio favouring own team).
    Near-uniform distribution. \\
  Layers 2--3 & Dead &
    Complete uniform distribution (7.69\,\% per token, 100\,\% entropy).
    No learned structure. \\
  \bottomrule
\end{tabularx}
\end{center}

\textbf{Effective model depth: 2~layers out of 4.}

\begin{keyinsight}
Post-norm transformers cause gradient vanishing in deeper layers.  The model
was effectively a 2-layer network paying minimal attention to the opponent.
This motivated the switch to a single layer with explicit attention bias.
\end{keyinsight}

% TODO: Insert attention heatmap visualisation showing:
% - Layer 0: structured pattern with some opponent attention
% - Layer 1: collapsed to near-uniform
% - Layers 2-3: completely uniform (dead)

\section{Pre-norm \& Architecture Experiments}

\textbf{Problem.} Attention collapse prevented the 4-layer transformer from
learning opponent-aware behaviour.  What architecture would actually work?

\textbf{Approach.} Systematic comparison of architectural variants, measuring
attention distribution and win rate:

\begin{center}
\small
\begin{tabularx}{\textwidth}{l c c c}
  \toprule
  \textbf{Configuration} &
  \textbf{CLS$\to$Opp Active} &
  \textbf{Our Active$\to$Opp Active} &
  \textbf{Win vs Heuristic} \\
  \midrule
  4L post-norm (baseline)    & 12.5\,\% & 3\,\%   & 0\,\% \\
  2L pre-norm               & 0\,\%    & ---      & 0\,\% \\
  2L pre-norm + attention bias & 12.5\,\% & ---   & 0\,\% \\
  \textbf{1L + attention bias} & \textbf{37.5\,\%} & \textbf{48.9\,\%} & \textbf{0\,\%} \\
  1L + bias + reward shaping & 62.5\,\% & ---      & 0\,\% \\
  \bottomrule
\end{tabularx}
\end{center}

\textbf{Key findings:}
\begin{itemize}
  \item Reducing from 4 to 2 layers with pre-norm \emph{regressed} opponent
        attention in layer~0 (from 12.5\,\% to 0\,\%).
  \item Adding attention bias restored and improved opponent attention.
  \item A single layer with attention bias produced the best result: the CLS
        token allocated 37.5\,\% of attention to the opponent's active
        Pok\'{e}mon, and the own active Pok\'{e}mon allocated 48.9\,\%.
  \item With reward shaping (matchup + action quality components), CLS
        attention to opponent active reached 62.5\,\%.
  \item \textbf{Despite improved attention, win rate against heuristic remained
        at 0\,\%} across all configurations.
\end{itemize}

\begin{keyinsight}
Architecture is no longer the bottleneck.  The model can \emph{see} the
opponent (48.9\,\% attention), but cannot translate that information into
winning actions.  The reward signal does not sufficiently incentivise
strategic play against competent opponents.
\end{keyinsight}

\section{Configuration Updates \& Loss Analysis}

Throughout the experimental phases, the model was scaled to approximately
10~million parameters.  A detailed analysis of training losses revealed a
consistent pattern:

\begin{center}
\small
\begin{tabularx}{0.8\textwidth}{l X}
  \toprule
  \textbf{Metric} & \textbf{Observation} \\
  \midrule
  Policy loss  & Very low throughout medium stage---the policy is not
                 making meaningful updates. \\
  Value loss   & Very high and does not decrease---the value network cannot
                 accurately predict returns. \\
  Entropy      & Remains high---the policy is uncertain and does not converge
                 to a confident strategy. \\
  Explained variance & Negative ($-0.003$)---the value function is worse than
                 predicting the mean. \\
  \bottomrule
\end{tabularx}
\end{center}

This loss signature is characteristic of a \textbf{value function failure}: the
critic cannot learn accurate state-value estimates, which impairs GAE advantage
estimation and consequently prevents the policy from improving.  The root cause
is likely the reward signal: with mixed-shaped components and sparse terminal
rewards, the value function faces a difficult credit-assignment problem.

% TODO: Insert MLflow loss curves showing the characteristic pattern:
% policy loss flat near zero, value loss high and noisy, entropy high.
% Annotate with the phase timeline (when each architectural change was made).
